% !PRINTED 198
\weberpage{198}
%<<CONT
\weq{6}{a_0 = 0, \quad a_1 = 1, \quad a_2 = 0, \quad a_3 = 0,}
then
\weq{7}{D = 0, \quad H = -x^2, \quad f = x^2 y, \quad Q = 2 x^3.}
Thus the left side would vanish while the right side does not
vanish, and in this lies a contradiction.

But finally we can also prove that on the left side of (4)
no power of 3 can occur which is not contained in all the
factors $M$, and which therefore cancels out.

In \S.\,2 we proved the theorem that the product of two
primitive integral rational functions of arbitrary variables is
again a primitive function. Here a primitive function is understood
to be one whose coefficients are integers without a common
divisor. The point now is to show that a sum of the form
\[
Q^\tau \, \Sigma \, M D^\gamma H^\alpha f^\beta,
\]
in which the $M$ are integers without a common divisor, cannot
become imprimitive, and in particular cannot acquire the factor 3,
when $Q$, $D$, $H$, $f$ are replaced by their expressions in the
$a$, $x$, $y$. By the theorem just mentioned it suffices, since $Q$
is a primitive function, to show this for the form
\weq{8}{\Sigma \, M D^\gamma H^\alpha f^\beta.}
In this we may moreover omit all terms whose $M$ already has
the factor 3, and finally, again by the theorem mentioned, we
may assume that the expression is not divisible by $D$, and
hence that it contains a term with $\gamma = 0$.
Suppose then that under these assumptions the expanded
expression (8) has the divisor 3, and now substitute the
particular values (6), (7), i.\,e.\ set $D = 0$, $H = -x^2$,
$f = x^2 y$; then the expression reduces to the single term
in which $\gamma = 0$, $2\alpha = \lambda$,
\[
(-1)^\alpha M x^{2(\alpha + \beta)} y^\beta
\]
and it would therefore follow that this $M$ must have the
divisor 3, which is contrary to the assumption. Thus we have
proved:

\gesperrt{Every integer covariant of the cubic form
can be expressed integrally and rationally and with integer coeffi-}
%CONT>>
